\section{Three Studies, Three Different Evaluation Units}
The research uses a public source cohort, a controlled-access external cohort and a separate negative-control dataset. Each was introduced for a distinct question. Study A asks whether a spectral classifier's answer changes across a five-minute recording. Study B builds and assesses the EEGraft decoder for AD versus controls. Study C freezes the two Study B models and measures how they behave on an external Normal-versus-broad-Dementia task. The Isfahan MCI investigation provides a separate example of apparently consistent but incorrect predictions. Its SVMs are not versions of the EEGraft decoder. The original participants in Studies A and B overlap in part, but the model protocols and partitions do not form a matched paired comparison.

The three main evaluation units must remain explicit: Study A counts people for folds but reports some stability states over correlated participant-repeat evaluations; Study B tests participants; Study C reports individual EEG \emph{recordings} because complete unique patient linkage is not available. Four-second windows, repeated folds and diagnostic classes are not additional independent patients.

\section{NEMAR/OpenNeuro Source and Analytical Subsets}
The NEMAR on004504 release mirrors OpenNeuro ds004504, originally described by Miltiadous and colleagues \cite{miltiadous2023dataset,miltiadous2026dataset}. It includes 88 resting, eyes-closed EEG participants: 36 AD, 23 FTD and 29 cognitively normal (CN). The source EEG has 19 channels acquired at 500 Hz, with recordings stored in EEGLAB \texttt{.set} format. Dataset-level \texttt{.json} metadata describes the collection; it is not a container for a multichannel waveform.

Study A retained 64 participants: 23 AD, 18 FTD and 23 controls. The AD/CN comparison has 46 people; FTD/CN has 41, reusing the same 23 controls. The eligibility record excludes 24 source participants for inadequate qualifying windows or review flags. Study B uses the full 65 AD/CN participants (36 AD, 29 CN) after excluding FTD from its binary task. It uses 52 development participants (29 AD, 23 CN) and a source locked holdout of 13 (7 AD, 6 CN). All 46 AD/CN IDs selected for Study A appear in the Study B manifest, but that overlap does not make their testing procedures the same.

\section{External CAUEEG Recording Cohort}
The Chung-Ang University Hospital EEG (CAUEEG) dataset introduced by Kim, Youn and Paik contains 1,379 clinical EEG recordings from 1,155 patients \cite{kim2023caueeg}. It includes 19 EEG channels, ECG and a photic channel; the sampled EEG signal rate is 200 Hz. Its official Dementia benchmark originally uses 1,187 records: 459 Normal, 417 mild cognitive impairment (MCI) and 311 Dementia. These three labels should not be read as 1,187 individually verified people, nor does the broad Dementia group mean that every record is confirmed Alzheimer's disease. The source provider distributes EDF files and annotation/event JSON through an application-based research-access process \cite{caueegofficial}.

The research uses the official \texttt{dementia-no-overlap} benchmark configuration, with the standard training split retained and known patient-overlap records removed from validation and test. The resulting partitions contain 1,122 \emph{recordings}: 950 training, 82 validation and 90 test. The primary binary task excludes MCI labels from the \emph{binary metrics and fitting}, not from the historical three-class benchmark counts. The binary sets contain 616 training recordings, 54 validation recordings and 57 test recordings. The 57 test recordings comprise 35 Normal and 22 Dementia; the remaining 33 test MCI records are descriptive only. Patient uniqueness within each split cannot be independently reconstructed from the restricted metadata available for this thesis.

\begin{table}[!htbp]\centering\small
\caption{CAUEEG official patient-overlap-reduced recording partitions, not unique participant counts}\label{tab:caueeg-splits}
\begin{tabular}{@{}lrrrrr@{}}\toprule
\textbf{Partition}&\textbf{Normal}&\textbf{MCI}&\textbf{Dementia}&\textbf{All}&\textbf{Binary N+D}\\\midrule
Training&367&334&249&950&616\\
Validation&34&28&20&82&54\\
Locked test&35&33&22&90&57\\\midrule
Total&436&395&291&1122&727\\\bottomrule
\end{tabular}\end{table}

\section{Cohort Selection and Source Evidence}
\TallThesisFigure[Cohort selection across the four analyses]{fig_04_01_cohort_flow.pdf}{Selection and evaluation pathways for NEMAR Studies A and B, external CAUEEG Study C, and the independent Isfahan negative control. The original 1,379 CAUEEG recordings are not the same as the 1,187 Dementia-task records or the 1,122 records retained by the official no-overlap benchmark. Units are explicitly marked; the diagram does not establish unique patient identities for every restricted recording.}{fig:cohort-flow}

Figure~\ref{fig:cohort-flow} distinguishes the four source populations. In particular, the 65 removed CAUEEG benchmark overlaps must not be confused with 192 recordings omitted from the original 1,379 because they were outside the three-class Dementia task. Full patient-by-split lineage requires restricted source linkage, so our uncertainty intervals describe recordings, not a fully verified patient-level sample.

\TallThesisFigure[Dataset metadata and available EEG evidence]{fig_04_02_dataset_evidence.pdf}{Dataset metadata and file-derived evidence. The NEMAR panel reproduces official metadata rather than a fabricated EEG trace. The Isfahan panel crops an uploaded six-second recording preview without asserting unverified physical amplitude units. The CAUEEG panel reports official split metadata without redistributing a restricted individual EEG trace. Signal and metadata preview evidence are deliberately distinguished.}{fig:dataset-evidence}

The dataset preview in Figure~\ref{fig:dataset-evidence} is a provenance illustration, not an example of model classification. For NEMAR, the publicly available metadata and raw \texttt{.set} file format are verifiable; no supplied record-specific NEMAR waveform was independently traced to an original sample interval. For Isfahan, the original display uses amplitude ticks that require source-level unit verification. For CAUEEG, we show counts rather than a restricted raw excerpt because display permission has not been established. No waveform was invented to fill either gap.

\section{Study A: Five-Minute Temporal Design}
The original five-minute SVM study extracts band-power and theta/alpha features from short spectral epochs in the 500 Hz NEMAR recordings. The archived spectral protocol describes two-second quality-controlled epochs, Welch spectral estimates and band ratios, followed by median aggregation. Its classifier is a standardized class-balanced radial-basis-function SVM with subject-disjoint nested cross-validation (five outer folds repeated five times and three inner folds). For each held-out subject, five consecutive 60-second EEG blocks are examined. First-minute balanced accuracy and subsequent changes or stable wrong outcomes are different measurements. Because the same people occur across repeats, the resulting 230 and 205 state records are correlated.

A separate duration investigation tested 30, 60, 120, 180 and 300 seconds at different recording positions; it is not the exact five-minute Study A protocol. Its archived row discrepancy was traced to an AD/CN file that contained both SVM and logistic-regression rows: 690 per model, 1,380 total. A saved source-to-source audit found no duplicate full prediction keys and agreed on first-minute labels after selecting the SVM records. The full repaired downstream duration recomputation and independent raw-EEG replay were not completed; therefore we do not infer an optimal recording duration from those archived plots.

\section{Study B: Source Decoder Development and Locked Test}
Study B originally held out 13 participants (7 AD, 6 CN) as a locked test. The remaining 52 (29 AD, 23 CN) were used for participant-disjoint development and final fitting. Five outer development folds cover 11, 11, 10, 10 and 10 participants. Each four-second window produces class probabilities, and participant scores aggregate over a maximum of 24 selected windows. Model selection is evaluated at the participant level. The fixed final DSRF and XGBoost models were assessed once on the 13-person source test. The 52 development participants and 13 test participants must not be blended to report development performance.

\section{Studies B and C: Input Preprocessing and Window Selection}
The executed source decoder adapter aligns 19 named EEG channels, converts source voltage units to microvolts, applies a fourth-order zero-phase 0.3--40 Hz Butterworth bandpass, resamples NEMAR data from 500 to 200 Hz, and then divides the microvolt signal by 100. CAUEEG already records at 200 Hz and retains 19 matched EEG channels from its 21 recorded channels; ECG and photic markers are not fed into the encoder. The 200 Hz sampling rate thus needs no CAUEEG resampling. The diagram documents the shared target format, not proof that every external recording has equivalent hardware filtering or reference conditions.

\TallThesisFigure[EEG preprocessing and window construction]{fig_04_03_preprocessing.png}{Implemented Study B/C decoder-input processing: channel alignment, 0.3--40 Hz passband, sampling alignment, fixed microvolt scaling, and up to 24 non-overlapping four-second windows sampled across each recording. Each input window contains 19 channels and 800 samples, arranged as four one-second patches of 200 samples. Signal sketches and filter response are illustrative, not recorded EEG. ICA and per-recording z-scoring are not implemented in the final decoder adapter.}{fig:preprocess}

For a long recording, up to 24 non-overlapping four-second windows are chosen across its duration by uniformly spaced indices. Thus the model receives at most 96 seconds of sampled signal, but these windows are not necessarily a consecutive 96-second excerpt. One model input is reshaped as $19\times4\times200$. The adapter does not add a new independent component analysis (ICA) step, common-average rereferencing, or per-recording z-score normalization. Original dataset derivatives may have their own processing history; those should not be confused with newly executed operations.

\section{Study C: Frozen Transfer and Assessment Boundaries}
The source-trained CBraMod--DSRF and XGBoost models are frozen and used for CAUEEG inference without fitting their weights to external labels. The original threshold 0.50 for both models is the primary external test comparison. An XGBoost threshold of 0.70 selected from CAUEEG training scores and separate train-fitted Platt maps are secondary exploratory variants. Their development checks occur on training and validation recordings; the locked test is not used to select the reported thresholds. The precise historical score-selection objective and whether the Platt fits used in-sample or out-of-fold scores were not reconstructable from the available execution records. These limits are stated again with the results.

\section{Independent Isfahan MCI Negative Control}
A separate historical MCI-versus-healthy EEG audit screened 61 EDF recordings. The quality gate retained 49: release EEG\_1 included 25 (11 MCI, 14 healthy) and EEG\_2 included 24 (12 MCI, 12 healthy). The two release-specific SVM evaluations were not pooled and were not produced by EEGraft's frozen decoder. The original identities across these two releases were not completely reconciled. This analysis is included solely to demonstrate why stable labels must be checked against correctness; it is not an external test of the same AD/CN model.

\section{Provenance and Reproducibility Boundary}
The original project contains NEMAR source code, manifests, saved protocol settings, model/checkpoint digests, development tables and source locked-test outcomes. The separate Ubuntu extension produced CAUEEG aggregate Code 19 results and Codes 21--25 audits, followed by Code 26 synthesis. These results were reported through terminal summaries and figure assets, but the original restricted CAUEEG JSON, EDF files, train/test prediction CSVs and model checkpoint binaries were not included in the reproducibility materials. The archived evidence supports numerical consistency checks against reported aggregate outputs but not a full independent rerun of the external models. The public report includes aggregate, nonidentifying evidence; unreproducible or confidential records are not invented or repackaged.
